\documentclass[12pt,a4paper]{article}
\usepackage[utf8]{inputenc}
\usepackage[indonesian]{babel}
\usepackage{amsmath}
\usepackage{amsfonts}
\usepackage{amssymb}
\usepackage{geometry}

\geometry{margin=2.5cm}

\title{Menghitung Garis Singgung dan Turunan Kedua Parametrik}
\author{Ahmad Fakhrul Bawani}
\date{}

\begin{document}

\maketitle

\section*{Soal}

Find an equation for the line tangent to the curve at the point defined by the given value of $t$. Also, find the value of $\frac{d^2y}{dx^2}$ at this point.
\begin{equation*}
  x = 16\sin t, \quad y = 4\cos t, \quad t = \frac{\pi}{4}
\end{equation*}

\subsection*{1. Menentukan Titik Singgung $(x_0, y_0)$}
Substitusikan nilai parameter $t = \frac{\pi}{4}$ ke dalam persamaan komponen posisi:
\begin{align*}
  x_0 &= 16\sin\left(\frac{\pi}{4}\right) = 16\left(\frac{1}{2}\sqrt{2}\right) = 8\sqrt{2} \\
  y_0 &= 4\cos\left(\frac{\pi}{4}\right) = 4\left(\frac{1}{2}\sqrt{2}\right) = 2\sqrt{2}
\end{align*}
Maka titik koordinat singgungnya berada pada \textbf{$\left(8\sqrt{2}, 2\sqrt{2}\right)$}.

\subsection*{2. Menghitung Turunan Pertama $\frac{dy}{dx}$ dan Gradien}
Gunakan aturan rantai untuk mencari komponen turunan terhadap parameter $t$:
\begin{align*}
  \frac{dx}{dt} &= \frac{d}{dt}(16\sin t) = 16\cos t \\
  \frac{dy}{dt} &= \frac{d}{dt}(4\cos t) = -4\sin t
\end{align*}

Maka fungsi rumus $\frac{dy}{dx}$ berbentuk:
\begin{equation*}
  \frac{dy}{dx} = \frac{\frac{dy}{dt}}{\frac{dx}{dt}} = \frac{-4\sin t}{16\cos t} = -\frac{1}{4}\tan t
\end{equation*}

Substitusikan $t = \frac{\pi}{4}$ untuk mencari gradien garis singgung ($m$):
\begin{equation*}
  m = \left. \frac{dy}{dx} \right|_{t=\frac{\pi}{4}} = -\frac{1}{4}\tan\left(\frac{\pi}{4}\right) = -\frac{1}{4}(1) = -\frac{1}{4}
\end{equation*}

\subsection*{3. Menyusun Persamaan Garis Singgung}
Menggunakan rumus persamaan garis $y - y_0 = m(x - x_0)$ dengan $m = -\frac{1}{4}$ dan titik singgung $\left(8\sqrt{2}, 2\sqrt{2}\right)$:
\begin{align*}
  y - 2\sqrt{2} &= -\frac{1}{4}\left(x - 8\sqrt{2}\right) \\
  y - 2\sqrt{2} &= -\frac{1}{4}x + 2\sqrt{2} \\
  y &= -\frac{1}{4}x + 4\sqrt{2}
\end{align*}

\subsection*{4. Menghitung Turunan Kedua $\frac{d^2y}{dx^2}$}
Rumus turunan kedua untuk persamaan berbentuk fungsi parametrik didefinisikan sebagai:
\begin{equation*}
  \frac{d^2y}{dx^2} = \frac{\frac{d}{dt}\left(\frac{dy}{dx}\right)}{\frac{dx}{dt}}
\end{equation*}

Mari cari komponen pembilang turunan dari hasil fungsi $\frac{dy}{dx} = -\frac{1}{4}\tan t$:
\begin{equation*}
  \frac{d}{dt}\left(-\frac{1}{4}\tan t\right) = -\frac{1}{4}\sec^2 t
\end{equation*}

Gabungkan kembali ke dalam rumus utama dengan $\frac{dx}{dt} = 16\cos t$:
\begin{equation*}
  \frac{d^2y}{dx^2} = \frac{-\frac{1}{4}\sec^2 t}{16\cos t} = -\frac{1}{64}\sec^3 t
\end{equation*}

Sekarang, substitusikan nilai $t = \frac{\pi}{4}$ ke dalam fungsi turunan kedua tersebut. Ingat bahwa $\sec\left(\frac{\pi}{4}\right) = \sqrt{2}$:
\begin{align*}
  \frac{d^2y}{dx^2} &= -\frac{1}{64} \left[\sec\left(\frac{\pi}{4}\right)\right]^3 \\
  &= -\frac{1}{64} (\sqrt{2})^3 \\
  &= -\frac{1}{64} (2\sqrt{2}) = -\frac{\sqrt{2}}{32}
\end{align*}

\subsection*{Kesimpulan Isian Jawaban}
Berdasarkan pengerjaan di atas, isian box jawaban yang tepat adalah:
\begin{itemize}
  \item \textbf{Turunan Pertama ($\frac{dy}{dx}$):} $-\frac{1}{4}\tan t$ \quad (Nilainya di titik tersebut adalah $-\frac{1}{4}$)
  \item \textbf{Persamaan Garis Singgung:} $y = -\frac{1}{4}x + 4\sqrt{2}$
  \item \textbf{Nilai Turunan Kedua ($\frac{d^2y}{dx^2}$):} $-\frac{\sqrt{2}}{32}$
\end{itemize}

\end{document}